% titel: Blatt 5, A8: Ionenkristall oder Molekülverbindung?

\subsection*{Frage}
Welche der Verbindungen $\mathrm{NaCl,\ SO_2,\ CsBr,\ BCl_3,\ PCl_3,\ HCl,\ KNO_3,\
NH_3,\ LiF,\ CH_4,\ PH_3}$ bilden Ionenkristalle, welche Moleküle?

\subsection*{Lösung}
\textbf{Ionenkristalle} (Metall--Nichtmetall bzw.\ Ionenpaar):
\begin{itemize}
  \item $\mathrm{NaCl}$ ($\mathrm{Na^+}$/$\mathrm{Cl^-}$)
  \item $\mathrm{CsBr}$ ($\mathrm{Cs^+}$/$\mathrm{Br^-}$)
  \item $\mathrm{KNO_3}$ ($\mathrm{K^+}$ und das Ionenradikal $\mathrm{NO_3^-}$)
  \item $\mathrm{LiF}$ ($\mathrm{Li^+}$/$\mathrm{F^-}$)
\end{itemize}
\textbf{Molekülverbindungen} (nur kovalente Bindungen zwischen Nichtmetallen):
$\mathrm{SO_2}$, $\mathrm{BCl_3}$, $\mathrm{PCl_3}$, $\mathrm{HCl}$ (polar,
kovalent), $\mathrm{NH_3}$, $\mathrm{CH_4}$, $\mathrm{PH_3}$.

\begin{ergebnis}
Ionenkristalle: $\mathrm{NaCl,\ CsBr,\ KNO_3,\ LiF}$. Moleküle:
$\mathrm{SO_2,\ BCl_3,\ PCl_3,\ HCl,\ NH_3,\ CH_4,\ PH_3}$.
\end{ergebnis}

\begin{kontrolle}
Alle vier Ionenkristalle enthalten ein Metall (Na, Cs, K, Li); die übrigen
bestehen nur aus Nichtmetallen. $\mathrm{KNO_3}$ ist der Trickfall: Salz aus $\mathrm{K^+}$ und dem komplexen
Nitrat-Anion. Das deckt sich mit der Regel Metall/Nichtmetall-Elektronenübertragung.
\end{kontrolle}

\begin{merksatz}
Salz = Metall + Nichtmetall (oder Ionen wie $\mathrm{NH_4^+}$, $\mathrm{NO_3^-}$);
Molekül = nur Nichtmetalle mit Elektronen-Paarbindung. Die EN-Differenz und die
Formelpartner entscheiden.
\end{merksatz}
